\SourceLeaf{333}%
\Article{32}
Endlich ergeben sich die Werthe der sieben Bestandtheile von \(H\) folgendermaassen:
\[
\begin{aligned}
1)\quad&-\frac{kk\cos T^2}{aacc}\mathord{.}ss\\
2)\quad&-\frac{kk\cos T^2}{12aa(1-ee)cc}\mathord{.}(1-2eess+ees^4)\\
3)\quad&+\frac{kk\sin T^2}{12aacc}\mathord{.}ss\\
4)\quad&+\frac{eekk\cos T^2}{4aa(1-ee)}\mathord{.}ss\\
5)\quad&-\frac{kk\sin T^2}{24aacc}\mathord{.}ss\\
6)\quad&+\frac{kk\cos T^2}{12aacc}\mathord{.}ss\\
7)\quad&+\frac{kk\cos T^2}{24aa(1-ee)cc}\mathord{.}(1+ss-2ees^4)
\end{aligned}
\]
Die Glieder 1 und 6 destruiren einander; die übrigen ergeben durch ihre Vereinigung
\[
H=\frac{kk\sin T^2}{24aacc}\mathord{.}ss-\frac{kk\cos T^2}{24aa(1-ee)}\mathord{.}(1-10eess)
\]
woraus, in Gemässheit von \(l=\lambda(1+Hrr)\) hervorgeht
\[
l=\lambda(1+\tfrac1{24}\tau\tau-\frac{1-ee}{24k^4}\mathord{.}(1-10eess)\beta\beta)\quad .\ .\ .\ .\ .\ .\quad(11)
\]

Die Formeln 9, 10, 11, welche die Auflösung unsrer Aufgabe in sich fassen, unterscheiden sich von den Formeln III, II, I, (Artt. 25, 24, 23) bloss darin, dass jene innerhalb der Parenthesen da \(\tau\) und \(\beta\) haben, wo in diesen \(t\) und \(b\) steht, was, wie man leicht sieht, in den Endresultaten nur Unterschiede fünfter Ordnung hervorbringt: da nun jene, wie aus ihrer Ableitung erhellet, bis zur fünften Ordnung ausschl. genau sind, so ist bewiesen, dass auch die nach der ersten Methode gefundenen Formeln I, II, III (Art. 23---25) dieselbe Genauigkeit besitzen.

\Article{33}
Zur numerischen Berechnung wird man die Formeln 9, 10, 11 lieber in folgende logarithmische Form bringen, bei welcher offenbar der Grad der Genauigkeit unverändert bleibt; \(M\) bezeichnet darin den Modulus des gewählten Logarithmensystems:
